\section{Введение (TODO переделать)}

\subsection{Актуальность работы}
В современной мировой экономике стартапы и предпринимательские инициативы являются основной точкой развития инноваций и научного прогресса, имеют исключительное значение в развитии экономики~\cite{startups-meaning}. Однако рынок венчурных инвестиций характеризуется крайне высокой степенью риска: доля провалившихся стартапов варьируется от $60\%$ до $90\%$ в зависимости от страны, временного периода и источников статистики, но общепринятой величиной, тем не менее, считается $90\%$~\cite{failurerate}. Иными словами, в долгосрочном периоде в среднем выживает лишь десятая часть от всех возникающих стартапов. В этих условиях инвесторы, бизнес-ангелы и сами предприниматели нуждаются в инструментах, позволяющих оценить риски и потенциал проекта на ранней стадии.

Эмпирический подход, основанный на анализе больших массивов данных о реальных стартапах, позволяет выявить объективные закономерности и статистически значимые факторы, связанные с успехом или провалом. Визуальная аналитика при этом делает результаты наглядными и понятными для широкого круга заинтересованных лиц.

Таким образом, актуальность работы обусловлена, во-первых, потребностью в объективном и подкрепленном реальными данными исследовании факторов успеха стартапов, и во-вторых, использованием статистических методов и визуальной аналитики, позволяющих выявить и наглядно представить ключевые закономерности.


\subsection{Цель}

Целью данной работы является проведение эмпирического исследования факторов успеха стартапов --- выпускников акселератора Y Combinator --- с использованием методов визуальной аналитики и статистического анализа, а также выявление и визуализация ключевых предикторов успеха или провала предпринимательских инициатив.

\subsection{Объект исследования}

Объектом исследования являются стартапы, прошедшие акселератор Y Combinator в период с 2005 по 2025 год, информация о которых представлена в открытых датасетах на платформе Kaggle.

\subsection{Предмет исследования}

Предметом исследования выступают характеристики стартапов (тип команды основателей, отрасль, сезон батча, лаконичность описания, количество тегов, география, стадия развития и другие), их связь со статусом компании (активна, неактивна, поглощена, вышла на IPO), а также динамика изменения этих характеристик во времени.

\subsection{Методы исследования}

\subsubsection{Методы предобработки данных}

\begin{itemize}
    \item Объединение нескольких датасетов с разной структурой и временными срезами.
    \item Нормализация форматов характеристик.
    \item Обработка пропущенных значений.
    \item Создание производных признаков.
\end{itemize}

\subsubsection{Разведочный анализ данных (EDA)}

\begin{itemize}
    \item Описательная статистика числовых и категориальных признаков.
    \item Анализ распределений и пропусков.
    \item Визуализация.
\end{itemize}

\subsubsection{Статистический анализ}

\begin{itemize}
    \item Сравнительный анализ групп успешных и неуспешных стартапов с контролем по когортам (по году батча).
    \item Критерий Манна-Уитни для проверки значимости различий между группами~\cite{mann-whitney}.
    \item Корреляционный анализ (коэффициент Пирсона) для оценки связи числовых признаков с успехом~\cite{pearson-corr}.
    \item Анализ долей успешных стартапов в разрезе категориальных признаков.
\end{itemize}

\subsubsection{Анализ динамики}

Сравнение данных двух временных срезов (2023 и 2025 годы) для выявления изменений в статусе и характеристиках стартапов.

\subsubsection{Инструменты визуализации}

Для построения графиков использовались библиотеки Python, такие как matplotlib и seaborn.


\subsection{Задачи}

\begin{enumerate}
    \item Произвести обзор литературы и существующих подходов в области анализа факторов успеха стартапов
    \item Осуществить сбор и предобработку данных из открытых источников, сформировать аналитические датасеты
    \item Произвести разведочный анализ собранных данных (EDA)
    \item Сформулировать и проверить статистические гипотезы о предикторах успеха с контролем по когортам для исключения ошибки выжившего
    \item Выполнить корреляционный анализ и анализ динамики изменений (2023--2025)
    \item Сформулировать выводы и рекомендации по ключевым факторам успеха или неудач
\end{enumerate}

\subsection{План работы}

План представлен в Таблице~\ref{tab:plan}

\begin{table}[htbp]
\centering
\small
\caption{Календарный план проекта (25.01.2026\,--\,01.04.2026)}
\label{tab:plan}
\begin{tabular}{@{}>{\raggedright\arraybackslash}p{3.2cm} p{2.3cm} p{6.8cm} p{3.5cm}@{}}
\toprule
\textbf{Этап} & \textbf{Неделя / Период} & \textbf{Основные задачи} & \textbf{Результат} \\
\midrule

\textbf{Этап 1: Подготовка данных и EDA} & Неделя 1 \newline 25.01--31.01 & Сбор данных (Kaggle); настройка окружения; слияние источников; обработка пропусков & Объединенные датасеты \\
\addlinespace[3pt]
-- & Неделя 2 \newline 01.02--07.02 & Создание производных признаков (бинарный статус, тип основателей); EDA: распределения, пропуски & Датасет A и B + ноутбук \\
\addlinespace[3pt]
-- & Неделя 3 \newline 08.02--14.02 & Анализ категориальных признаков; визуализация батчей, отраслей, стран & EDA-отчет с графиками \\
\midrule

\textbf{Этап 2: Анализ предикторов} & Неделя 4 \newline 15.02--21.02 & Проверка гипотез H1--H4: соло vs команда, B2B vs Consumer, сезон батча, лаконичность описания & Графики и статистические тесты \\
\addlinespace[3pt]
-- & Неделя 5 \newline 22.02--28.02 & Проверка гипотез H5--H8: количество тегов, найм, география, стадия развития & Полный набор визуализаций \\
\addlinespace[3pt]
-- & Неделя 6 \newline 01.03--07.03 & Корреляционный анализ; анализ динамики 2023--2025; матрица переходов & Корреляционная матрица + выводы \\
\midrule

\textbf{Этап 3: Оформление} & Неделя 7 \newline 08.03--14.03 & Написание отчета: введение, описание данных, методология & Черновик курсовой работы \\
\addlinespace[3pt]
-- & Неделя 8 \newline 15.03--21.03 & Написание отчета: результаты, выводы, оформление графиков & Полный текст работы \\
\addlinespace[3pt]
-- & Финал \newline 22.03--01.04 & Проверка оформления; создание презентации; финализация & Готовый проект к защите \\
\bottomrule
\end{tabular}
\end{table}
